\section{Resultados y discusión}
\subsection{Parte A: simulación}
Se recibieron y revisaron los módulos \texttt{divisor}, \texttt{contador} y \texttt{top\_A}, el archivo de restricciones y el testbench \texttt{tb\_parte\_A}. Los cuatro archivos VHDL se analizaron con Vivado 2025.1 y se ejecutó el testbench con XSim. La simulación utilizó un reloj de 10 ns, \texttt{MAX=4} y un contador de tres bits para observar el desborde en un intervalo breve.

Durante 36 ciclos se compararon mediante \texttt{assert} el valor de \texttt{en} con un modelo de periodicidad y el valor del contador con una cuenta de referencia. Se verificaron el pulso de un ciclo cada cuatro ciclos, la retención entre habilitaciones, el incremento, el desborde de siete a cero y el reset inicial y durante la cuenta. La ejecución finalizó a los 385,001 ns con el mensaje de comprobaciones correctas, sin errores de aserción. En ese intervalo se consumieron ocho habilitaciones; el noveno pulso se generó al final del bucle y fue seguido por reset. El contador consume el enable registrado en el flanco anterior, por lo que su actualización ocurre un ciclo después de generarse el pulso.

Esta prueba ejercitó divisor y contador; el top de diez bits fue analizado, pero no fue la unidad instanciada por el testbench. La simulación no acredita funcionamiento físico ni valida el mapeo de pines del shield.
\completar{Incorporar captura legible de las formas de onda de esta ejecución.}

\subsection{Parte A: funcionamiento físico}
\completar{Registrar los valores observados durante al menos diez segundos, la conmutación de LED0, el incremento cada segundo y el comportamiento del reset. Si no se implementó, indicar expresamente la limitación.}

\subsection{Parte B: simulación}
\completar{Insertar captura de al menos dos ciclos de la FSM; identificar clk, rst, en, op, A, B, resultado y carry. Analizar reset y sumas con acarreo.}
\begin{figure}[!t]\centering
\fbox{\begin{minipage}[c][25mm][c]{0.85\columnwidth}\centering\small Captura de la simulación de la Parte B\end{minipage}}
\caption{Transiciones de la FSM y salidas del datapath. Imagen pendiente de incorporación.}\label{fig:simB}
\end{figure}

\subsection{Parte B: validación física}
\completar{Describir valores A y B usados, secuencia visible en LED9:8, resultado en LED3:0, carry en LED4 y reset. Comparar lo esperado con lo observado, sin marcar como verificados los casos aún no ensayados.}

\begin{table}[!t]
\centering\caption{Registro de pruebas de la FSM y del datapath}\label{tab:pruebas}
\begin{tabular}{@{}ccccc@{}}\toprule
\textbf{A} & \textbf{B} & \textbf{op} & \textbf{Esperado} & \textbf{Observado}\\\midrule
0101 & 0011 & 00 & 00001 & --\\
0101 & 0011 & 01 & 00111 & --\\
0101 & 0011 & 10 & 00110 & --\\
0101 & 0011 & 11 & 01000 & --\\
1111 & 0011 & 11 & 10010 & --\\\bottomrule
\end{tabular}
\end{table}
La Tabla~\ref{tab:pruebas} presenta casos de referencia calculados, no mediciones. El campo observado quedó sin completar hasta su contrastación experimental. Los cinco bits de salida se representaron como \texttt{LED4:0}.

\subsection{Discusión y limitaciones}
\completar{Analizar diferencias entre simulación y placa, problemas concretos de restricciones, síntesis, lectura de switches o temporización; explicar sus causas y las soluciones efectivamente aplicadas. Valorar la legibilidad del video e indicar si las tomas permitieron verificar las entradas y salidas.}

